\documentclass[11pt]{article}
\usepackage[letterpaper,margin=1in]{geometry}
\usepackage[T1]{fontenc}
\usepackage[utf8]{inputenc}
\usepackage{microtype}
\usepackage[hidelinks=false,colorlinks=true,urlcolor=blue]{hyperref}
\setlength{\parindent}{0pt}
\setlength{\parskip}{6pt}
\pagestyle{empty}

\begin{document}

[Date]

AI's 10 to Watch Selection Committee\\
IEEE Intelligent Systems

Dear Members of the Selection Committee,

I am pleased to strongly recommend Rahul Gupta for IEEE Intelligent Systems' AI's 10 to Watch. My recommendation is grounded in Rahul's sustained contributions to responsible AI, trustworthy machine learning, fairness, privacy, red-teaming, evaluation, and the translation of AI research into deployed industrial systems.

I know Rahul through the responsible-AI and trustworthy-AI research community, where his work stands out for combining technical substance with real deployment impact. He has built a rare profile for an early-career AI researcher: he publishes in major AI venues, invents patented technologies, leads responsible-AI programs inside one of the world's most consequential AI organizations, and creates mechanisms that support academic researchers working on trustworthy AI.

Rahul's scientific contributions are substantial. His work includes BOLD, a widely used dataset and metric framework for measuring bias in open-ended language generation; fairness methods for text classification and large-scale natural-language-understanding systems; FLIRT, a feedback-loop approach for in-context red teaming; privacy-preserving NLP methods; model-unlearning and data-curation work; rigorous agentic benchmark design; and recent first-author work on CBRN uplift evaluation for frontier language models. These contributions are directly aligned with the central questions facing modern AI: how to measure, mitigate, and govern harms in systems that are increasingly powerful, open-ended, and embedded in real products.

What distinguishes Rahul is that he has not treated responsible AI as only a publication area. At Amazon, he founded and led Responsible AI science for Alexa, later led Responsible AI for Amazon Nova, and now serves as Responsible AI lead for Amazon's AGI / Frontier AI Assets organization. In Alexa, his work helped translate fairness and privacy research into mechanisms for improving experiences for underserved customer populations. In Nova and AGI, he has led work across safety evaluations, red-teaming, fairness, privacy-sensitive behavior, mitigation, external assessment, and release-readiness mechanisms for frontier foundation models.

Rahul's patent portfolio is an important part of this translation from research to practice. He is an inventor of 25+ patented or patent-pending technologies. His Alexa-era inventions include mechanisms related to cohort determination, model configuration, user-data processing, spoken-language understanding, and NLP component evaluation. These technologies support more equitable, fair, and private AI experiences in production settings. His Nova-era patent filings, including work on generative language models and synthetic-data generation, support responsible-AI methods that improve the safety profile of frontier foundation models.

His impact also extends to the academic ecosystem. Rahul has helped organize TrustNLP, Trustworthy Speech Processing, and the SemEval Unlearning Challenge, and he has delivered invited keynotes, talks, and panels on responsible AI and frontier-model evaluation. He also led the Amazon Nova AI Challenge, which attracted more than 100 university applications and supported 10 university teams working on trusted AI through substantial research sponsorships, AWS credits, mentoring, and prize opportunities. Across these efforts, Rahul has helped create venues and resources through which academic researchers can test ideas, receive recognition, and advance safety, fairness, privacy, unlearning, red-teaming, and robustness research.

AI's 10 to Watch should recognize researchers whose work is already significant and whose future influence is likely to grow. Rahul clearly meets that standard. He has produced influential research artifacts, helped deploy responsible-AI mechanisms in large-scale systems, and built bridges between industry and academia at a moment when the AI community urgently needs exactly that kind of leadership.

I recommend Rahul Gupta with enthusiasm.

Sincerely,

Kush R. Varshney\\
IBM Fellow, IBM Research\\
Email / signature placeholder

\end{document}
